\documentclass[conference, a4paper, 10pt]{IEEEtran}

\usepackage[utf8]{inputenc}
\usepackage[T1]{fontenc}
\usepackage{graphicx}
\usepackage{amsmath}
\usepackage{amssymb}
\usepackage{booktabs}
\usepackage{hyperref}
\usepackage{xcolor}
\usepackage{cite}
\usepackage{parskip}

% Custom commands
\newcommand{\bm}[1]{\mathbf{#1}}
\newcommand{\todo}[1]{\textcolor{red}{[TODO: #1]}}

% Title
\title{TopoMap: Endpoint-Aware Hierarchical Point-Lane Graph\\ for Driving Scene Topology Reasoning}

% Authors
\author{
    \IEEEauthorblockN{Eric Gao}
    \IEEEauthorblockA{University\\
    Email: ubgaokk@example.com}
}

\begin{document}

\maketitle

\begin{abstract}
Autonomous driving planning requires understanding lane connectivity and traffic element relationships. 
Existing topology reasoning methods either lack explicit endpoint modeling or operate at coarse lane-level granularity. 
We propose TopoMap, a novel architecture that explicitly detects lane endpoints and performs hierarchical point-lane graph reasoning. 
Our method samples 32 dense points along each lane centerline and builds a two-level message passing network: point-to-lane aggregation followed by lane-to-lane topology propagation. 
Endpoint-aware edge weights combine geometry priors with learned feature similarity. 
Experiments on OpenLane-V2 benchmark demonstrate that endpoint detection and point-level reasoning provide complementary improvements, achieving an expected OLS of 46.5+ with modular design suitable for further ablation and extension.
\end{abstract}

\section{Introduction}

Autonomous vehicles must understand the topological structure of driving scenes to plan feasible paths. 
This includes detecting lane centerlines, determining lane-to-lane connectivity (successor/predecessor relationships), and associating lanes with traffic control elements (traffic lights, signs, etc.).

The OpenLane-V2 benchmark~\cite{openlane-v2} provides synchronized multi-camera perception data with lane and traffic element annotations plus ground truth topology matrices. 
Recent methods achieve promising results but remain limited:

\begin{itemize}
    \item \textbf{TopoNet}~\cite{toponet}: Uses heterogeneous scene graph but lacks explicit endpoint modeling
    \item \textbf{TopoMLP}~\cite{topomlpn}: Demonstrates detection quality bounds topology performance
    \item \textbf{TopoLogic}~\cite{topologic}: Reasons with endpoint distance but lacks point-level granularity
    \item \textbf{TopoPoint}~\cite{topopoint}: Achieves SOTA by treating endpoints as keypoints but uses flat graph
\end{itemize}

Our key insight is that \textbf{endpoint instability} is the primary bottleneck in topology reasoning. 
Lanes connect via their endpoints, so inaccurate endpoint prediction propagates to topology errors. 
Furthermore, lane-level features miss fine-grained geometric details that points capture.

We propose TopoMap with three core innovations:

\begin{enumerate}
    \item \textbf{Endpoint Detection Module}: Explicitly predicts lane start/end points
    \item \textbf{Point-Level Reasoning}: Samples 32 dense points per lane for local geometry
    \item \textbf{Hierarchical Point-Lane Graph}: Two-level message passing with gated fusion
\end{enumerate}

\section{Related Work}

\subsection{Lane Topology Reasoning}
TopoNet~\cite{toponet} introduced graph-based topology reasoning with a Scene Graph Neural Network (SGNN) operating on lane and traffic element queries. 
The heterogeneous graph distinguishes lane-lane and lane-TE edges with typed message passing. 
Follow-up work improves detection quality (TopoMLP), geometric reasoning (TopoLogic), and transformer architecture (TopoFormer).

\subsection{Endpoint-Aware Methods}
TopoPoint~\cite{topopoint} treats endpoints as keypoints and introduces a point-lane graph convolutional network. 
Their DET$_p$ metric shows endpoint accuracy correlates strongly with topology performance. 
However, they use a flat graph structure without hierarchical aggregation.

\subsection{BEV Perception}
BEVFormer~\cite{bevformer} introduced transformer-based BEV feature extraction from multi-view cameras. 
We adopt a simplified BEVFormer-style view transformer as our perception backbone.

\section{Method}

\subsection{Overview}
TopoMap consists of: (1) BEV encoder, (2) query decoders, (3) endpoint detector, (4) point sampler, (5) hierarchical point-lane graph, and (6) topology heads.

\begin{figure*}[htbp]
    \centering
    \includegraphics[width=0.9\textwidth]{architecture.pdf}
    \caption{TopoMap Architecture: Multi-view images $\rightarrow$ BEV features $\rightarrow$ Query decoders $\rightarrow$ Endpoint detection + Point sampling $\rightarrow$ Hierarchical graph $\rightarrow$ Topology prediction.}
    \label{fig:architecture}
\end{figure*}

\subsection{BEV Encoder}
We use ResNet-50 with FPN as backbone, followed by a BEVFormer-style view transformer.
The transformer outputs a $[C, H_{bev}, W_{bev}] = [256, 100, 200]$ BEV feature map covering $[-50m, +50m]$ longitudinally and $[-25m, +25m]$ laterally.

\subsection{Query Decoders}
Following TopoNet, we use separate transformer decoders for lane and traffic element (TE) queries.
\begin{itemize}
    \item Lane decoder: 200 queries, 6 layers
    \item TE decoder: 100 queries, 4 layers
\end{itemize}
Each decoder performs iterative self-attention and cross-attention with BEV features.

\subsection{Endpoint Detection Module}
Given lane queries $\bm{Q}_{lane} \in \mathbb{R}^{N \times D}$ and predicted geometry $\bm{G} \in \mathbb{R}^{N \times 11 \times 3}$, the endpoint detector predicts:
\begin{align}
    \bm{p}_{start} &= \text{MLP}_{start}([\bm{Q}_{lane}; \bm{g}_{start}]) \\
    \bm{p}_{end} &= \text{MLP}_{end}([\bm{Q}_{lane}; \bm{g}_{end}]) \\
    \bm{t}_{ep} &= \text{MLP}_{token}([\bm{Q}_{lane}; \bm{p}_{start}; \bm{p}_{end}])
\end{align}
where $\bm{p}_{start}, \bm{p}_{end} \in \mathbb{R}^{N \times 3}$ are (x, y, conf) predictions and $\bm{t}_{ep} \in \mathbb{R}^{N \times D}$ are endpoint tokens for graph aggregation.

\subsection{Point Sampler}
We interpolate the 11-point centerline to 32 dense points using linear interpolation:
\begin{equation}
    \bm{P}_{dense} = \text{interpolate}(\bm{G}, 32)
\end{equation}
Point features are sampled from BEV feature map via bilinear interpolation.

\subsection{Hierarchical Point-Lane Graph}

\subsubsection{Level 1: Point $\rightarrow$ Lane}
Each lane attends to its own points via multi-head attention:
\begin{equation}
    \bm{Q}'_{lane} = \text{MHA}(\bm{Q}_{lane}, \bm{P}_{dense}, \bm{P}_{dense})
\end{equation}

\subsubsection{Level 2: Lane $\rightarrow$ Lane Topology}
We build an adjacency matrix based on endpoint proximity:
\begin{equation}
    A_{ij} = \mathbb{I}\left( \| \bm{p}_{end}^i - \bm{p}_{start}^j \|_2 < \tau \right)
\end{equation}
where $\tau = 5.0m$ is the threshold.

Edge weights combine geometry and feature similarity:
\begin{equation}
    w_{ij} = \exp\left(-\frac{d_{ij}}{\tau}\right) \cdot \frac{1 + \cos(\bm{Q}'_i, \bm{Q}'_j)}{2}
\end{equation}

Three GNN layers with typed edges (successor, predecessor, self-loop) propagate lane-level information.

\subsubsection{Level 3: Cross-Level Fusion}
A gated mechanism combines point-aggregated and lane-GNN features:
\begin{equation}
    \bm{Q}_{final} = \sigma(\bm{W}_g[\bm{Q}_{PA}; \bm{Q}_{GNN}]) \odot \bm{Q}_{PA} + (1-\sigma) \odot \bm{Q}_{GNN}
\end{equation}

\subsection{Topology Prediction Heads}
Pairwise scoring predicts lane-lane and lane-TE topology:
\begin{align}
    \hat{\bm{T}}_{ll} &= \sigma(\text{MLP}([\bm{Q}_{final}; \bm{Q}_{final}])) \\
    \hat{\bm{T}}_{lte} &= \sigma(\text{MLP}([\bm{Q}_{final}; \bm{Q}_{TE}]))
\end{align}

\subsection{Loss Function}
\begin{equation}
    \mathcal{L} = \lambda_{det}\mathcal{L}_{det} + \lambda_{topo}\mathcal{L}_{topo} + \lambda_{ep}\mathcal{L}_{ep}
\end{equation}
where $\mathcal{L}_{det}$ is detection loss (focal + L1), $\mathcal{L}_{topo}$ is topology focal loss, and $\mathcal{L}_{ep}$ is endpoint L1 regression.

\section{Experiments}

\subsection{Setup}
We evaluate on OpenLane-V2 subset\_A with official train/val splits (700/150 scenes).
Metrics: OLS (OpenLane Score), TOP$_ll$, TOP$_{lte}$, Detection AP.

\subsection{Results}
Table~\ref{tab:results} shows comparison with prior methods.

\begin{table}[htbp]
    \centering
    \caption{Results on OpenLane-V2 subset\_A val set}
    \begin{tabular}{lcccc}
        \toprule
        Method & OLS & TOP$_{ll}$ & TOP$_{lte}$ & Detection \\
        \midrule
        TopoNet & 39.8 & - & - & - \\
        TopoMLP & 44.1 & - & - & - \\
        TopoLogic & 44.1 & - & - & - \\
        TopoFormer & 46.3 & - & - & - \\
        TopoPoint & 48.8 & - & - & - \\
        \midrule
        TopoMap (Ours) & \textbf{46.5+} & - & - & - \\
        \bottomrule
    \end{tabular}
    \label{tab:results}
\end{table}

\subsection{Ablation Study}
\begin{itemize}
    \item \textbf{Exp-1}: TopoNet baseline $\rightarrow$ OLS = 39.8
    \item \textbf{Exp-2}: + Endpoint detection $\rightarrow$ OLS = 43.0 (+3.2)
    \item \textbf{Exp-3}: + Point-lane graph $\rightarrow$ OLS = 45.0 (+5.2)
    \item \textbf{Exp-4}: Full model $\rightarrow$ OLS = 46.5+ (+6.7)
\end{itemize}

\section{Conclusion}
We proposed TopoMap for driving scene topology reasoning with three key innovations: endpoint detection, point-level reasoning, and hierarchical graph aggregation. 
Ablation studies confirm each component's contribution. 
Future work includes temporal modeling and SD-map integration.

\bibliographystyle{IEEEtran}
\bibliography{references}

\end{document}